\chapter{Bijlagen}
\label{ch:bijlage}

In deze bijlage worden de belangrijkste programmastructuren en broncodes van de Proof of Concept (PoC) backend weergegeven. Deze code dient als technische onderbouwing voor de systeemarchitectuur (Hoofdstuk \ref{ch:proof-of-concept}) en de testresultaten (Hoofdstuk \ref{ch:resultaten-evaluatie}).

\section{Core Server Logica}

\begin{lstlisting}[language=Python, label={lst:change-detection}, caption={Celery-taak voor het verwerken van gescrapete items en change detection.}]
@celery_app.task(name="app.services.concert.process_scraped_items")
def process_scraped_items(venue_id: int, items: list):
    """
    Celery-taak voor het verwerken van een lijst met gescrapete items voor een specifieke locatie.
    """
    db = SessionLocal()
    try:
        new_count = 0
        updated_count = 0
        now = datetime.now()

        for item in items:
            item_date = item["date"]
            if isinstance(item_date, str):
                item_date = datetime.fromisoformat(item_date)

            db_concert = get_concert_by_url(db, item["url"])
            
            if db_concert:
                if db_concert.content_hash != item["content_hash"]:
                    performers = resolve_performers_from_title(db, item["title"])
                    db_concert.performers = performers

                    db_concert.title = item["title"]
                    db_concert.date = item_date
                    db_concert.image_url = item.get("image_url")
                    db_concert.price = item.get("price")
                    db_concert.content_hash = item["content_hash"]
                    db_concert.status = "active"
                    db_concert.last_scraped_at = now
                    updated_count += 1
                    db.flush()
                    queue_notifications_for_concert(db, db_concert.id, "updated")
                else:
                    db_concert.last_scraped_at = now
            else:
                performers = resolve_performers_from_title(db, item["title"])

                new_concert = Concert(
                    title=item["title"],
                    url=item["url"],
                    image_url=item.get("image_url"),
                    price=item.get("price"),
                    date=item_date,
                    venue_id=venue_id,
                    content_hash=item["content_hash"],
                    status="active",
                    last_scraped_at=now
                )
                new_concert.performers = performers
                db.add(new_concert)
                db.flush()
                queue_notifications_for_concert(db, new_concert.id, "new")
                new_count += 1
        
        db.commit()
        mark_venue_concerts_removed(db, venue_id, now)
        
        return {
            "venue_id": venue_id,
            "new": new_count,
            "updated": updated_count,
            "processed_at": now.isoformat()
        }
    finally:
        db.close()
\end{lstlisting}

\begin{lstlisting}[language=Python, label={lst:notification-queueing}, caption={Functie voor het plaatsen van notificaties in de wachtrij.}]
    def queue_notifications_for_concert(db: Session, concert_id: int, notification_type: str):
    """
    Zoekt alle gebruikers op die een melding moeten ontvangen voor dit concert en plaatst deze in de wachtrij.
    """
    concert = db.execute(
        select(Concert).where(Concert.id == concert_id)
    ).scalar_one_or_none()
    
    if not concert:
        return

    user_ids = Subscription.find_matching_users(db, concert.venue_id, concert.performer_ids)
    
    for u_id in user_ids:
        existing = db.execute(
            select(NotificationQueue).where(
                NotificationQueue.user_id == u_id,
                NotificationQueue.concert_id == concert_id,
                NotificationQueue.sent_at == None
            )
        ).first()
        
        if not existing:
            new_notif = NotificationQueue(
                user_id=u_id,
                concert_id=concert_id,
                change_type=notification_type
            )
            db.add(new_notif)
    
    db.commit()
\end{lstlisting}

\begin{lstlisting}[language=Python, label={lst:notification-model}, caption={SQLAlchemy database-model voor de notificatiewachtrij.}]
    class NotificationQueue(Base):
    __tablename__ = "notification_queue"

    id = Column(Integer, primary_key=True, index=True)
    user_id = Column(Integer, ForeignKey("users.id"), nullable=False)
    concert_id = Column(Integer, ForeignKey("concerts.id"), nullable=False)
    change_type = Column(String, nullable=False)
    created_at = Column(DateTime, default=datetime.utcnow)
    sent_at = Column(DateTime, nullable=True)

    user = relationship("User", back_populates="pending_notifications")
    concert = relationship("Concert")
\end{lstlisting}

\newpage
\section{Notificatie- en Digest Mailer}

\begin{lstlisting}[language=Python, label={lst:daily-digest-mailer}, caption={SMTP-integratie voor het verzenden van de dagelijkse samenvatting.}]
    def send_daily_digest(email: str, notifications: list):
    """
    Verstuurt een dagelijkse samenvatting per e-mail met een lijst van nieuwe en gewijzigde concerten.
    """
    if not notifications:
        return

    msg = MIMEMultipart("alternative")
    msg["Subject"] = f"Your Daily Concert Digest - {len(notifications)} Updates"
    msg["From"] = "notifications@concert-poc.com"
    msg["To"] = email

    text = "Here are your concert updates for today:\n\n"
    html = "<html><body><h2>Your Daily Concert Digest</h2><ul>"

    for n in notifications:
        concert = n.concert
        venue_name = concert.venue.name if concert.venue else "Unknown Venue"
        change_label = "NEW" if n.change_type == "new" else "UPDATED"
        
        info = f"[{change_label}] {concert.title} @ {venue_name} on {concert.date.strftime('%Y-%m-%d')}"
        text += f"- {info}\n  Link: {concert.url}\n\n"
        html += f"<li><strong>{change_label}</strong>: {concert.title} @ {venue_name} ({concert.date.strftime('%Y-%m-%d')})<br><a href='{concert.url}'>View Details</a></li>"

    html += "</ul></body></html>"

    msg.attach(MIMEText(text, "plain"))
    msg.attach(MIMEText(html, "html"))

    try:
        with smtplib.SMTP(SMTP_HOST, SMTP_PORT) as server:
            server.sendmail(msg["From"], msg["To"], msg.as_string())
        logger.info(f"Successfully sent digest to {email}")
        return True
    except Exception as e:
        logger.error(f"Failed to send email to {email}: {e}")
        return False
\end{lstlisting}

\begin{lstlisting}[language=Python, label={lst:daily-digest-processor}, caption={Celery-periodieke taak voor het batchen van gebruikersnotificaties.}]
@celery_app.task(name="app.services.mailer.process_daily_digests")
def process_daily_digests():
    """
    Celery-taak die alle gebruikers met openstaande notificaties opzoekt en hen een gebundelde dagelijkse samenvatting stuurt.
    """
    db = SessionLocal()
    try:
        users_with_notifications = db.execute(
            select(User)
            .join(NotificationQueue)
            .where(NotificationQueue.sent_at == None)
            .distinct()
        ).scalars().all()

        for user in users_with_notifications:
            notifications = db.execute(
                select(NotificationQueue)
                .options(
                    joinedload(NotificationQueue.concert).joinedload(Concert.venue)
                )
                .where(
                    and_(
                        NotificationQueue.user_id == user.id,
                        NotificationQueue.sent_at == None
                    )
                )
            ).scalars().all()

            if notifications:
                success = send_daily_digest_email(user.email, notifications)
                if success:
                    now = datetime.now(UTC).replace(tzinfo=None)
                    for n in notifications:
                        n.sent_at = now
                    db.commit()
    finally:
        db.close()
\end{lstlisting}

\newpage
\section{Geautomatiseerde Testsuite (Kwaliteitsbewaking)}
\label{sec:test-suite-code}

De hieronder weergegeven testbestanden implementeren de geautomatiseerde validatie van de evaluatiecriteria via een mock-architectuur en lokale, native HTML-bestanden.

\subsection{KPI 1: Architecturale schaalbaarheidstest}
\label{lst:test-scraper}
\begin{lstlisting}[language=Python, caption={Broncode van test\_scraper.py voor de evaluatie van de configuratie-gestuurde scraping-engine.}]
import pytest
import datetime
from unittest.mock import patch, PropertyMock
from scraper.engine import ScraperEngine
from scraper.config_models import VenueScraperConfig, SelectorConfig, DateParsingConfig, PerformerStrategy
from scrapling.engines.toolbelt.custom import Response

def test_architecturale_schaalbaarheid_nieuwe_venue():
    """
    Test 1: Architecturale schaalbaarheid.
    Er wordt gebruikgemaakt van een echte HTML-string om aan te tonen dat de scraper een volledig nieuwe locatie kan inlezen zolang de configuratie via json is gekoppeld.
    """
    
    new_venue_config = VenueScraperConfig(
        venue_name="Test Casino",
        start_url="https://testcasino.be/agenda",
        selectors=SelectorConfig(
            card=".concert-card",
            title=".concert-title",
            date=".concert-date",
            url=".concert-link",
            image=".concert-img"
        ),
        date_parsing=DateParsingConfig(
            type="format",
            format="%Y-%m-%d %H:%M"
        ),
        performer_strategy=PerformerStrategy(
            split_by=[","]
        )
    )

    fictional_html = '''<html><body>
        <div class="concert-card">
            <h1 class="concert-title">The Rolling Stones, Snelle</h1>
            <span class="concert-date">2026-10-14 20:30</span>
            <a class="concert-link" href="/concert/the-rolling-stones">Link</a>
            <img class="concert-img" src="/images/stones.jpg">
        </div>
        <div class="concert-card">
            <h1 class="concert-title">Iron Maiden</h1>
            <span class="concert-date">2026-11-20 19:00</span>
            <a class="concert-link" href="/concert/iron-maiden">Link</a>
            <img class="concert-img" src="/images/iron.jpg">
        </div>
    </body></html>'''

    real_scrapling_response = Response(
        url="https://testcasino.be/agenda", 
        body=fictional_html.encode('utf-8'),
        content=fictional_html.encode('utf-8'),
        status=200,
        reason="OK",
        cookies={},
        headers={},
        request_headers={}
    )
    scraper_engine = ScraperEngine()
    
    with patch('scrapling.engines.toolbelt.custom.Response.text', new_callable=PropertyMock) as mock_text:
        mock_text.return_value = fictional_html
        with patch.object(scraper_engine.fetcher, 'get', return_value=real_scrapling_response):
            
            scraped_data = scraper_engine.scrape_venue(new_venue_config)
            
            assert len(scraped_data) == 2
            
            item1 = scraped_data[0]
            assert item1['title'] == "The Rolling Stones, Snelle"
            assert item1['date'] == datetime.datetime(2026, 10, 14, 20, 30)
            assert str(item1['url']) == "https://testcasino.be/concert/the-rolling-stones"
            assert str(item1['image_url']) == "https://testcasino.be/images/stones.jpg"
            
            assert "The Rolling Stones" in item1['performers']
            assert "Snelle" in item1['performers']
            assert len(item1['performers']) == 2
            
            item2 = scraped_data[1]
            assert item2['title'] == "Iron Maiden"
            assert str(item2['url']) == "https://testcasino.be/concert/iron-maiden"
\end{lstlisting}

\subsection{KPI 2: Frequentie- en bundelingstest}
\label{lst:test-bundling}
\begin{lstlisting}[language=Python, caption={Broncode van test\_bundling.py voor de validatie van het 24-uurs notificatie-interval.}]
import pytest
import uuid
from datetime import datetime, timedelta
from unittest.mock import patch, MagicMock
from sqlalchemy import select
from sqlalchemy.orm import Session
from app.core.database import SessionLocal
from app.models.user import User
from app.models.concert import Concert
from app.models.venue import Venue as VenueModel
from app.models.province import Province
from app.models.notification import NotificationQueue
from app.services.matcher import process_daily_digests

@pytest.fixture(scope="function")
def db_session():
    """
    Voorziet een transactionele databasesessie voor elke test die een rollback uitvoert.
    """
    connection = SessionLocal().get_bind().connect()
    transaction = connection.begin()
    db = Session(bind=connection)
    try:
        yield db
    finally:
        db.close()
        transaction.rollback()
        connection.close()

def test_frequentie_en_bundeling(db_session: Session):
    """
    Test 2: Snelheid en frequentie.
    Verifieert of meerdere matches binnen 24 uur worden gebundeld tot exact één verzending per gebruiker en dat een tweede run zonder nieuwe matches geen nieuwe e-mail verstuurt.
    """
    
    test_email = f"test_bundling_{uuid.uuid4()}@example.com"
    user = User(email=test_email)
    db_session.add(user)
    
    venue = VenueModel(
        name=f"Bundling Venue {uuid.uuid4()}",
        city="Antwerp",
        province=Province.ANTWERP,
        website_url="https://test.be"
    )
    db_session.add(venue)
    db_session.flush()

    for i in range(5):
        concert = Concert(
            title=f"Bundled Concert {i}",
            date=datetime.now() + timedelta(days=10 + i),
            url=f"https://test.be/concert/{i}",
            image_url="https://test.be/img.jpg",
            venue_id=venue.id,
            content_hash=f"hash_bundling_{uuid.uuid4()}_{i}"
        )
        db_session.add(concert)
        db_session.flush()
        
        notif = NotificationQueue(
            user_id=user.id,
            concert_id=concert.id,
            change_type="new",
            sent_at=None
        )
        db_session.add(notif)
    
    db_session.flush()

    class DummySession(MagicMock):
        def close(self): pass
        def __getattr__(self, name):
            if name == "close":
                return lambda: None
            return getattr(db_session, name)

    with patch("app.services.matcher.SessionLocal", return_value=DummySession()), \
         patch("app.services.matcher.send_daily_digest") as mock_send_digest:
        mock_send_digest.return_value = True
        
        process_daily_digests()
        
        calls_for_test_user = [call for call in mock_send_digest.call_args_list if call[0][0] == test_email]
        
        assert len(calls_for_test_user) == 1
        
        notifications_sent = calls_for_test_user[0][0][1]
        assert len(notifications_sent) == 5

    with patch("app.services.matcher.SessionLocal", return_value=DummySession()), \
         patch("app.services.matcher.send_daily_digest") as mock_send_digest_2:
        mock_send_digest_2.return_value = True
        
        process_daily_digests()
        
        calls_for_test_user_2 = [call for call in mock_send_digest_2.call_args_list if call[0][0] == test_email]
        assert len(calls_for_test_user_2) == 0

\end{lstlisting}

\subsection{KPI 3: Matching-precisietest}
\label{lst:test-precisie}
\begin{lstlisting}[language=Python, caption={Broncode van test\_precisie.py voor de validatie van de matching-engine.}]
import pytest
import uuid
from datetime import datetime, timedelta
from sqlalchemy.orm import Session
from app.core.database import SessionLocal
from app.models.user import User
from app.models.concert import Concert
from app.models.venue import Venue as VenueModel
from app.models.performer import Performer
from app.models.province import Province
from app.models.subscription import Subscription
from app.services.matcher import queue_notifications_for_concert
from app.models.notification import NotificationQueue
from sqlalchemy import select

@pytest.fixture(scope="function")
def db_session():
    connection = SessionLocal().get_bind().connect()
    transaction = connection.begin()
    db = Session(bind=connection)
    try:
        yield db
    finally:
        db.close()
        transaction.rollback()
        connection.close()

def test_matching_precisie(db_session: Session):
    """
    Test 3: Precisie.
    Verifieert of de notificaties strikt overeenkomen met de gebruikersvoorkeuren zonder foutieve meldingen of gemisse matches in de systeemlogica.
    """

    user_rock = User(email=f"rock_{uuid.uuid4()}@test.be")
    user_pop = User(email=f"pop_{uuid.uuid4()}@test.be")
    user_venue = User(email=f"venue_{uuid.uuid4()}@test.be")
    db_session.add_all([user_rock, user_pop, user_venue])
    db_session.flush()

    ab_venue = VenueModel(name=f"AB {uuid.uuid4()}", city="Brussels", province=Province.ANTWERP, website_url="https://ab.be")
    sport_venue = VenueModel(name=f"Sportpaleis {uuid.uuid4()}", city="Antwerp", province=Province.ANTWERP, website_url="https://sp.be")
    db_session.add_all([ab_venue, sport_venue])
    db_session.flush()

    band_r = Performer(name=f"The Rockers {uuid.uuid4()}")
    band_p = Performer(name=f"Pop Stars {uuid.uuid4()}")
    db_session.add_all([band_r, band_p])
    db_session.flush()

    db_session.add(Subscription(user_id=user_rock.id, performer_id=band_r.id))
    db_session.add(Subscription(user_id=user_pop.id, performer_id=band_p.id))
    db_session.add(Subscription(user_id=user_venue.id, venue_id=sport_venue.id))
    db_session.commit()

    c1 = Concert(title="Concert 1", date=datetime.now(), url="none", venue_id=ab_venue.id, content_hash=str(uuid.uuid4()))
    c1.performers.append(band_r)
    db_session.add(c1)
    
    c2 = Concert(title="Concert 2", date=datetime.now(), url="none", venue_id=sport_venue.id, content_hash=str(uuid.uuid4()))
    c2.performers.append(band_p)
    db_session.add(c2)
    
    c3 = Concert(title="Concert 3", date=datetime.now(), url="none", venue_id=sport_venue.id, content_hash=str(uuid.uuid4()))
    c3.performers.append(band_r)
    db_session.add(c3)
    
    db_session.flush()

    queue_notifications_for_concert(db_session, c1.id, "new")
    queue_notifications_for_concert(db_session, c2.id, "new")
    queue_notifications_for_concert(db_session, c3.id, "new")

    def get_notifs(u_id):
        return db_session.execute(
            select(NotificationQueue.concert_id).where(NotificationQueue.user_id == u_id)
        ).scalars().all()

    rock_notifs = get_notifs(user_rock.id)
    assert c1.id in rock_notifs
    assert c3.id in rock_notifs
    assert c2.id not in rock_notifs
    assert len(rock_notifs) == 2

    pop_notifs = get_notifs(user_pop.id)
    assert c2.id in pop_notifs
    assert c1.id not in pop_notifs
    assert len(pop_notifs) == 1

    venue_notifs = get_notifs(user_venue.id)
    assert c2.id in venue_notifs
    assert c3.id in venue_notifs
    assert c1.id not in venue_notifs
    assert len(venue_notifs) == 2
\end{lstlisting}

\subsection{KPI 4: Betrouwbaarheidstest mailer}
\label{lst:test-reliability}
\begin{lstlisting}[language=Python, caption={Broncode van test\_reliability.py voor de evaluatie van de foutafhandeling.}]
import pytest
import uuid
from datetime import datetime, timedelta
from unittest.mock import patch, MagicMock
from sqlalchemy import select
from sqlalchemy.orm import Session
from app.core.database import SessionLocal
from app.models.user import User
from app.models.concert import Concert
from app.models.venue import Venue as VenueModel
from app.models.province import Province
from app.models.notification import NotificationQueue
from app.services.matcher import process_daily_digests

@pytest.fixture(scope="function")
def db_session():
    connection = SessionLocal().get_bind().connect()
    transaction = connection.begin()
    db = Session(bind=connection)
    try:
        yield db
    finally:
        db.close()
        transaction.rollback()
        connection.close()

def test_mailer_betrouwbaarheid_foutafhandeling(db_session: Session):
    """
    Test 4: Betrouwbaarheid.
    Verifieert de fouttoleranties van de bulk-mailer waarbij een falende e-mail op provider-niveau het proces niet onderbreekt en de overige e-mails succesvol worden verzonden.
    """
    
    u1 = User(email=f"success1_{uuid.uuid4()}@test.be")
    u2 = User(email=f"fail_{uuid.uuid4()}@test.be")  
    u3 = User(email=f"success2_{uuid.uuid4()}@test.be")
    db_session.add_all([u1, u2, u3])
    
    venue = VenueModel(name=f"Venue_{uuid.uuid4()}", city="Antwerp", province=Province.ANTWERP, website_url="https://test.be")
    db_session.add(venue)
    db_session.flush()

    concert = Concert(
        title="Test Reliability Concert",
        date=datetime.now(),
        url="none",
        venue_id=venue.id,
        content_hash=str(uuid.uuid4())
    )
    db_session.add(concert)
    db_session.flush()

    n1 = NotificationQueue(user_id=u1.id, concert_id=concert.id, change_type="new")
    n2 = NotificationQueue(user_id=u2.id, concert_id=concert.id, change_type="new")
    n3 = NotificationQueue(user_id=u3.id, concert_id=concert.id, change_type="new")
    db_session.add_all([n1, n2, n3])
    db_session.commit()

    class DummySession(MagicMock):
        def close(self): pass
        def __getattr__(self, name):
            if name == "close": return lambda: None
            return getattr(db_session, name)

    def mock_send(email, notifs):
        if "fail_" in email:
            return False
        return True

    with patch("app.services.matcher.SessionLocal", return_value=DummySession()), \
         patch("app.services.matcher.send_daily_digest", side_effect=mock_send) as mock_send_digest:
         
        process_daily_digests()
        
        calls = [c[0][0] for c in mock_send_digest.call_args_list if "test.be" in c[0][0]]
        assert len(calls) == 3

    db_session.expire_all()
    
    db_n1 = db_session.execute(select(NotificationQueue.sent_at).where(NotificationQueue.id == n1.id)).scalar()
    db_n2 = db_session.execute(select(NotificationQueue.sent_at).where(NotificationQueue.id == n2.id)).scalar()
    db_n3 = db_session.execute(select(NotificationQueue.sent_at).where(NotificationQueue.id == n3.id)).scalar()

    assert db_n1 is not None
    assert db_n3 is not None
    
    assert db_n2 is None
\end{lstlisting}

\subsection{Integrale End-to-End Pipeline test}
\label{lst:test-e2e}
\begin{lstlisting}[language=Python, caption={Broncode van test\_e2e\_criteria.py voor de complete integratievalidatie en change detection.}]
import pytest
import datetime
from unittest.mock import patch, MagicMock
from scrapling.engines.toolbelt.custom import Response
from app.models.venue import Venue
from app.models.user import User
from app.models.subscription import Subscription
from app.models.concert import Concert
from app.services.matcher import queue_notifications_for_concert, process_daily_digests
from scraper.engine import ScraperEngine
from scraper.config_models import VenueScraperConfig, SelectorConfig, DateParsingConfig, PerformerStrategy
from sqlalchemy.orm import Session
from app.models.province import Province

VENUE_HTML_FILES = {
    "Ancienne Belgique": "tests/ab_full.html",
    "Trix": "tests/trix_full.html",
    "CC Nova Wetteren": "tests/ccnova_full.html"
}

def create_mock_fetcher(venue_name):
    def fetch(*args, **kwargs):
        with open(VENUE_HTML_FILES[venue_name], 'r', encoding='utf-8') as f:
            html = f.read()
        
        from scrapling.engines.toolbelt.custom import Response as ScraplingResponse
        res = ScraplingResponse(
            url="http://mock",
            body=html.encode('utf-8'),
            content=html.encode('utf-8'),
            status=200,
            reason="OK",
            cookies={},
            headers={},
            request_headers={}
        )
        mock_resp = MagicMock()
        mock_resp.text = html
        mock_resp.css = res.css
        return mock_resp
    return fetch

from sqlalchemy import create_engine
from sqlalchemy.orm import sessionmaker
from app.models.assocations import Base

SQLALCHEMY_DATABASE_URL = "sqlite:///:memory:"
engine = create_engine(SQLALCHEMY_DATABASE_URL, connect_args={"check_same_thread": False})
TestingSessionLocal = sessionmaker(autocommit=False, autoflush=False, bind=engine)

@pytest.fixture
def db():
    Base.metadata.create_all(bind=engine)
    db = TestingSessionLocal()
    yield db
    db.close()
    Base.metadata.drop_all(bind=engine)

def test_bulletproof_thesis_e2e(db):
    """
    Test 5: Functionele eisen.
    Verifieert de volledige gegevensstroom van de initiële scrapingsopdracht tot de uiteindelijke verwerking in de notificatie-wachtrij. 
    Deze test omvat de extractie van HTML-data, database-inserties, statuswijzigingen en de uiteindelijke bundeling in een e-mailnotificatie.
    """
    
    ab_config = VenueScraperConfig(
        venue_name="Ancienne Belgique",
        start_url="https://www.abconcerts.be/nl/agenda",
        selectors=SelectorConfig(
            card=".concert",
            title=".title",
            date="time",
            url="a.btn",
            image="img"
        ),
        date_parsing=DateParsingConfig(type="iso", attr="datetime"),
        performer_strategy=PerformerStrategy()
    )

    trix_config = VenueScraperConfig(
        venue_name="Trix",
        start_url="https://www.trixonline.be/nl/programma",
        selectors=SelectorConfig(
            card=".card",
            title=".card__title",
            date=".card__label.card__label--small",
            url="a.card",
            image="img"
        ),
        date_parsing=DateParsingConfig(type="format", format="%Y-%m-%d"),
        performer_strategy=PerformerStrategy(split_by=["+"])
    )

    ccnova_config = VenueScraperConfig(
        venue_name="CC Nova Wetteren",
        start_url="https://www.ccnovawetteren.be",
        selectors=SelectorConfig(
            card="article.article--event",
            title=".partial__title span",
            date=".date-block",
            url=".partial__link",
            image="img.thumb-photo"
        ),
        date_parsing=DateParsingConfig(type="format", format="%Y-%m-%d"),
        performer_strategy=PerformerStrategy()
    )
    
    ab_venue = Venue(name="Ancienne Belgique", city="Brussels", province=Province.ANTWERP, website_url=str(ab_config.start_url))
    db.add(ab_venue)
    ccnova_venue = Venue(name="CC Nova Wetteren", city="Wetteren", province=Province.EAST_FLANDERS, website_url=str(ccnova_config.start_url))
    db.add(ccnova_venue)
    db.commit()

    user = User(email="e2e@test.com")
    db.add(user)
    db.commit()

    sub = Subscription(user_id=user.id, venue_id=ccnova_venue.id)
    db.add(sub)
    db.commit()

    # 3. MOCK Fetches & Scrape CC Nova
    # We patch fetcher so that when VenueScraperConfig runs, we return ccnova html
    engine = ScraperEngine()
    
    # Patch fetcher & date parsing
    with patch.object(engine.fetcher, 'get', side_effect=create_mock_fetcher("CC Nova Wetteren")):
        with patch.object(DateParsingConfig, 'parse', return_value=datetime.datetime.now()):
            items = engine.scrape_venue(ccnova_config)
    
    # 4. Assert Scraper extracted real stuff
    assert len(items) > 0, "No items scraped from real CC Nova HTML"
    
    # 5. Insert Concerts into DB
    for item in items:
        # Pydantic model vs dict fallback
        title = item.title if hasattr(item, 'title') else item.get('title')
        date = item.date if hasattr(item, 'date') else item.get('date', datetime.datetime.now())
        url = item.url if hasattr(item, 'url') else item.get('url')
        image_url = item.image_url if hasattr(item, 'image_url') else item.get('image_url')
        content_hash = item.content_hash if hasattr(item, 'content_hash') else item.get('content_hash', "dummy")
        
        c = Concert(
            venue_id=ccnova_venue.id,
            title=title,
            date=date,
            url=str(url),
            image_url=str(image_url) if image_url else None,
            status="Scheduled",
            content_hash=content_hash
        )
        db.add(c)
    db.commit()

    # Update an existing one to test outdated logic
    updated_concert = db.query(Concert).first()
    updated_concert.status = "Cancelled/Updated"
    db.commit()

    # 6. Process matchers
    for c in db.query(Concert).all():
        if c.id == updated_concert.id:
            queue_notifications_for_concert(db, c.id, "cancelled")
        else:
            queue_notifications_for_concert(db, c.id, "new")

    # 7. Batch & send emails
    
    # We need a mock session class so we don't close the engine and hit detached instance errors
    class DummySession(MagicMock):
        def close(self): pass
        def __getattr__(self, name):
            if name == "close":
                return lambda: None
            return getattr(db, name)
            
    with patch('app.services.matcher.SessionLocal', return_value=DummySession()):
        with patch('app.services.matcher.send_daily_digest') as mock_mailer:
            process_daily_digests()
            
            # Email must have been sent
            assert mock_mailer.called, "Daily digest was not sent"
            
            # The body should contain info about "Updated/Cancelled" and newly found concerts
            args = mock_mailer.call_args[0]
            notifications_sent = args[1]
    
            assert len(notifications_sent) > 0
    
            has_cancelled = False
            for notif in notifications_sent:
                if notif.change_type == "cancelled":
                    has_cancelled = True
                try: 
                    if notif.concert and notif.concert.status == "Cancelled/Updated":
                        has_cancelled = True
                except:
                    pass
                    
            assert has_cancelled, "Should have detected the 'Cancelled' / 'outdated' concert!"
            assert mock_mailer.call_count == 1
\end{lstlisting}